% GENERATED FILE. Edit canonical lessons, then run npm run generate:books.
% canonical-source-hash: fnv1a64:f049da308bb153bb
% canonical-lessons: ML-C377-bhukampam, ML-C377-urulpottal, ML-C377-oliccupokuka, ML-C377-orrappetuka, ML-C377-kanatakuka

\chapter{Earthquake and Landslide}
\label{ch:ml-earthquake-and-landslide}
\hlchaptermodality{377}

\begin{chapteropening}
\textbf{By the end of this chapter:} \emph{I can say an earthquake, a landslide, to be washed away, to be cut off and to go missing.}

Complete the last lesson of chapter 377: I can say an earthquake, a landslide, to be washed away, to be cut off and to go missing.
\end{chapteropening}

\section[\texorpdfstring{\ml{ഭൂകമ്പം} (bhūkampaṁ)}{ഭൂകമ്പം (bhūkampaṁ)}]{\ml{ഭൂകമ്പം} (bhūkampaṁ) — an earthquake}

\label{lesson:ML-C377-bhukampam}



\begin{quote}
Before the new one: say the Malayalam for to get hurt, then the Malayalam for to drown.
\end{quote}

\subsection*{You'll want to know: \ml{ഭൂകമ്പം}}
\textbf{\ml{ഭൂകമ്പം}} — \emph{bhūkampaṁ} — \textquotedblleft{}an earthquake\textquotedblright{}.

\begin{quote}
\ml{ഭൂകമ്പത്തിൽ} \ml{നഗരത്തിലെ} \ml{ഒരുപാട്} \ml{കെട്ടിടങ്ങൾ} \ml{തകർന്നു}.
\end{quote}

\begin{quote}
\emph{bhūkampattil nagarattile orupāṭŭ keṭṭiṭaṅṅaḷ takarnnu.}
\end{quote}

\begin{quote}
Many buildings in the city collapsed in the earthquake.
\end{quote}

\ml{ഭൂ}, the earth, + \ml{കമ്പം}, trembling.

\begin{cousinweb}
\textbf{\ml{ഭൂകമ്പം}} (\emph{bhūkampaṁ}) is borrowed from Sanskrit \emph{bhūkampa}, \textquotedblleft{}an earthquake\textquotedblright{} (CDIAL 9547), made of \emph{bhū}, \textquotedblleft{}the earth\textquotedblright{}, and \emph{kampa}, \textquotedblleft{}trembling\textquotedblright{}.
\end{cousinweb}

\begin{grammarlens}[title={What you've built}]
One more word for telling what happened.
\end{grammarlens}

\subsection*{Your turn}
\begin{itemize}
\raggedright
  \item \emph{Say it:} \emph{bhūkampaṁ}
  \item \emph{Say it:} \emph{bhūkampaṁ}, once more
  \item \emph{Say it:} \emph{bhūkampattil nagarattile orupāṭŭ keṭṭiṭaṅṅaḷ takarnnu.} — Many buildings in the city collapsed in the earthquake.
  \item \emph{Recall:} say the Malayalam for to get hurt, then the Malayalam for to drown, then say \emph{bhūkampaṁ} again
\end{itemize}

\begin{tcolorbox}[breakable,colback=teal!4,colframe=teal!35!black,title={Before you move on}]
What does \ml{ഭൂകമ്പം} mean? (\textquotedblleft{}An earthquake\textquotedblright{}.) Say it once more.
\end{tcolorbox}
\section[\texorpdfstring{\ml{ഉരുൾപൊട്ടൽ} (uruḷpoṭṭal)}{ഉരുൾപൊട്ടൽ (uruḷpoṭṭal)}]{\ml{ഉരുൾപൊട്ടൽ} (uruḷpoṭṭal) — a landslide}

\label{lesson:ML-C377-urulpottal}



\begin{quote}
Before the new one: say the Malayalam for to drown, then the Malayalam for an earthquake.
\end{quote}

\subsection*{You'll want to know: \ml{ഉരുൾപൊട്ടൽ}}
\textbf{\ml{ഉരുൾപൊട്ടൽ}} — \emph{uruḷpoṭṭal} — \textquotedblleft{}a landslide\textquotedblright{}.

\begin{quote}
\ml{മലയിൽ} \ml{ഉരുൾപൊട്ടൽ} \ml{ഉണ്ടായി}.
\end{quote}

\begin{quote}
\emph{malayil uruḷpoṭṭal uṇṭāyi.}
\end{quote}

\begin{quote}
There was a landslide on the hill.
\end{quote}

Common on the hills in the rainy season.

\begin{cousinweb}
\textbf{\ml{ഉരുൾപൊട്ടൽ}} (\emph{uruḷpoṭṭal}) is a native Dravidian compound of \ml{ഉരുൾ}, \textquotedblleft{}a rolling\textquotedblright{} (DEDR 664), and \ml{പൊട്ടൽ}, \textquotedblleft{}a bursting\textquotedblright{}, from \ml{പൊട്ടുക} (DEDR 4490): the hillside bursts and rolls down.
\end{cousinweb}

\begin{grammarlens}[title={What you've built}]
One more word for telling what happened.
\end{grammarlens}

\subsection*{Your turn}
\begin{itemize}
\raggedright
  \item \emph{Say it:} \emph{uruḷpoṭṭal}
  \item \emph{Say it:} \emph{uruḷpoṭṭal}, once more
  \item \emph{Say it:} \emph{malayil uruḷpoṭṭal uṇṭāyi.} — There was a landslide on the hill.
  \item \emph{Recall:} say the Malayalam for to drown, then the Malayalam for an earthquake, then say \emph{uruḷpoṭṭal} again
\end{itemize}

\begin{tcolorbox}[breakable,colback=teal!4,colframe=teal!35!black,title={Before you move on}]
What does \ml{ഉരുൾപൊട്ടൽ} mean? (\textquotedblleft{}A landslide\textquotedblright{}.) Say it once more.
\end{tcolorbox}
\section[\texorpdfstring{\ml{ഒലിച്ചുപോകുക} (oliccupōkuka)}{ഒലിച്ചുപോകുക (oliccupōkuka)}]{\ml{ഒലിച്ചുപോകുക} (oliccupōkuka) — to be washed away}

\label{lesson:ML-C377-oliccupokuka}



\begin{quote}
Before the new one: say the Malayalam for an earthquake, then the Malayalam for a landslide.
\end{quote}

\subsection*{You'll want to know: \ml{ഒലിച്ചുപോകുക}}
\textbf{\ml{ഒലിച്ചുപോകുക}} — \emph{oliccupōkuka} — \textquotedblleft{}to be washed away\textquotedblright{}.

\begin{quote}
\ml{വെള്ളപ്പൊക്കത്തിൽ} \ml{പാലം} \ml{ഒലിച്ചുപോയി}.
\end{quote}

\begin{quote}
\emph{veḷḷappokkattil pālaṁ oliccupōyi.}
\end{quote}

\begin{quote}
The bridge was washed away in the flood.
\end{quote}

Past: \ml{ഒലിച്ചുപോയി}, \emph{oliccupōyi}.

\ml{ഒലിക്കുക}, to flow, + \ml{പോകുക}, to go.

\begin{cousinweb}
\textbf{\ml{ഒലിച്ചുപോകുക}} (\emph{oliccupōkuka}) joins the native verb \ml{ഒലിക്കുക}, \textquotedblleft{}to flow, to run as water\textquotedblright{} (DEDR 999), and \ml{പോകുക} (\emph{pōkuka}), \textquotedblleft{}to go\textquotedblright{}: to go off with the flowing water.
\end{cousinweb}

\begin{grammarlens}[title={What you've built}]
One more word for telling what happened.
\end{grammarlens}

\subsection*{Your turn}
\begin{itemize}
\raggedright
  \item \emph{Say it:} \emph{oliccupōkuka}
  \item \emph{Say it:} \emph{oliccupōkuka}, once more
  \item \emph{Say it:} \emph{veḷḷappokkattil pālaṁ oliccupōyi.} — The bridge was washed away in the flood.
  \item \emph{Recall:} say the Malayalam for an earthquake, then the Malayalam for a landslide, then say \emph{oliccupōkuka} again
\end{itemize}

\begin{tcolorbox}[breakable,colback=teal!4,colframe=teal!35!black,title={Before you move on}]
What does \ml{ഒലിച്ചുപോകുക} mean? (\textquotedblleft{}To be washed away\textquotedblright{}.) Say it once more.
\end{tcolorbox}
\section[\texorpdfstring{\ml{ഒറ്റപ്പെടുക} (oṟṟappeṭuka)}{ഒറ്റപ്പെടുക (oṟṟappeṭuka)}]{\ml{ഒറ്റപ്പെടുക} (oṟṟappeṭuka) — to be cut off; to be isolated}

\label{lesson:ML-C377-orrappetuka}



\begin{quote}
Before the new one: say the Malayalam for a landslide, then the Malayalam for to be washed away.
\end{quote}

\subsection*{You'll want to know: \ml{ഒറ്റപ്പെടുക}}
\textbf{\ml{ഒറ്റപ്പെടുക}} — \emph{oṟṟappeṭuka} — \textquotedblleft{}to be cut off; to be isolated\textquotedblright{}.

\begin{quote}
\ml{മഴ} \ml{കാരണം} \ml{ഗ്രാമം} \ml{ഒറ്റപ്പെട്ടു}.
\end{quote}

\begin{quote}
\emph{maḻa kāraṇaṁ grāmaṁ oṟṟappeṭṭu.}
\end{quote}

\begin{quote}
The village was cut off because of the rain.
\end{quote}

Past: \ml{ഒറ്റപ്പെട്ടു}, \emph{oṟṟappeṭṭu}.

From \ml{ഒറ്റ}, single, alone: also of a lonely person.

\begin{cousinweb}
\textbf{\ml{ഒറ്റപ്പെടുക}} (\emph{oṟṟappeṭuka}) is native Dravidian: \ml{ഒറ്റ} (\emph{oṟṟa}), \textquotedblleft{}single, alone\textquotedblright{}, which DEDR 990 sets with the word for one, and -\ml{പ്പെടുക}, \textquotedblleft{}to come to be\textquotedblright{}.
\end{cousinweb}

\begin{grammarlens}[title={What you've built}]
One more word for telling what happened.
\end{grammarlens}

\subsection*{Your turn}
\begin{itemize}
\raggedright
  \item \emph{Say it:} \emph{oṟṟappeṭuka}
  \item \emph{Say it:} \emph{oṟṟappeṭuka}, once more
  \item \emph{Say it:} \emph{maḻa kāraṇaṁ grāmaṁ oṟṟappeṭṭu.} — The village was cut off because of the rain.
  \item \emph{Recall:} say the Malayalam for a landslide, then the Malayalam for to be washed away, then say \emph{oṟṟappeṭuka} again
\end{itemize}

\begin{tcolorbox}[breakable,colback=teal!4,colframe=teal!35!black,title={Before you move on}]
What does \ml{ഒറ്റപ്പെടുക} mean? (\textquotedblleft{}To be cut off; to be isolated\textquotedblright{}.) Say it once more.
\end{tcolorbox}
\section[\texorpdfstring{\ml{കാണാതാകുക} (kāṇātākuka)}{കാണാതാകുക (kāṇātākuka)}]{\ml{കാണാതാകുക} (kāṇātākuka) — to go missing}

\label{lesson:ML-C377-kanatakuka}



\begin{quote}
Before the new one: say the Malayalam for to be washed away, then the Malayalam for to be cut off.
\end{quote}

\subsection*{You'll want to know: \ml{കാണാതാകുക}}
\textbf{\ml{കാണാതാകുക}} — \emph{kāṇātākuka} — \textquotedblleft{}to go missing\textquotedblright{}.

\begin{quote}
\ml{മൂന്ന്} \ml{ദിവസം} \ml{മുൻപ്} \ml{ഒരു} \ml{കുട്ടിയെ} \ml{കാണാതായി}.
\end{quote}

\begin{quote}
\emph{mūnnŭ divasaṁ munpŭ oru kuṭṭiye kāṇātāyi.}
\end{quote}

\begin{quote}
A child went missing three days ago.
\end{quote}

Past: \ml{കാണാതായി}, \emph{kāṇātāyi}.

The missing person takes -\ml{എ}: \ml{അവനെ} \ml{കാണാതായി}, literally it became not seeing him.

\begin{cousinweb}
\textbf{\ml{കാണാതാകുക}} (\emph{kāṇātākuka}) is native Dravidian: the negative of \ml{കാണുക} (\emph{kāṇuka}), \textquotedblleft{}to see\textquotedblright{} (DEDR 1443), with \ml{ആകുക}, \textquotedblleft{}to become\textquotedblright{}: someone comes to be not seen.
\end{cousinweb}

\begin{grammarlens}[title={What you've built}]
One more word for telling what happened.
\end{grammarlens}

\subsection*{Your turn}
\begin{itemize}
\raggedright
  \item \emph{Say it:} \emph{kāṇātākuka}
  \item \emph{Say it:} \emph{kāṇātākuka}, once more
  \item \emph{Say it:} \emph{mūnnŭ divasaṁ munpŭ oru kuṭṭiye kāṇātāyi.} — A child went missing three days ago.
  \item \emph{Recall:} say the Malayalam for to be washed away, then the Malayalam for to be cut off, then say \emph{kāṇātākuka} again
\end{itemize}

\begin{tcolorbox}[breakable,colback=teal!4,colframe=teal!35!black,title={Before you move on}]
What does \ml{കാണാതാകുക} mean? (\textquotedblleft{}To go missing\textquotedblright{}.) Say it once more.
\end{tcolorbox}
